\documentclass[11pt]{article}

\usepackage[utf8]{inputenc}
\usepackage[T1]{fontenc}
\usepackage{booktabs}
\usepackage{amsmath}
\usepackage{graphicx}
\usepackage{siunitx}
\usepackage[margin=1in]{geometry}
\usepackage{hyperref}

\sisetup{detect-all = true}
\graphicspath{{figures/}}

% This is Chapter 5: number equations, tables, and figures as 5.x to match the prose
% references (Equation 5.1, Table 5.1, Figure 5.1) and the Markdown version.
\renewcommand{\theequation}{5.\arabic{equation}}
\renewcommand{\thetable}{5.\arabic{table}}
\renewcommand{\thefigure}{5.\arabic{figure}}
\renewcommand{\thesection}{5.\arabic{section}}

% Give tables room to breathe: taller rows and wider column gaps.
\renewcommand{\arraystretch}{1.35}
\setlength{\tabcolsep}{8pt}

% Shorthands for the dataset names so underscores never reach text mode.
\newcommand{\Dd}{\ensuremath{D_d}}
\newcommand{\Dr}{\ensuremath{D_r}}

\title{Chapter 5: Results and Analysis}
\author{RepairA11y}
\date{}

\begin{document}
\maketitle

\begin{quote}
\textbf{Note on data provenance.} Every number in this chapter is produced by an
evaluation script in this repository and is traceable to a saved result file under
\texttt{experiments/<name>/results/}. The exact command and result file for each
figure are listed in the companion \texttt{README.md}. The numbers (the RQ1
ablation, the safety check measured on the same runs, and the \Dr{} detection and
live repair) come from recorded live runs dated 2026-06-20 and 2026-06-21. The
recorded runs are the source of truth because re-running the full LLM suite requires
the OpenRouter free-tier quota, which was exhausted at the time of writing. Where a
quantity could not be verified from a script it is tagged \texttt{[VERIFY]} with an
explanation of what is needed to confirm it.
\end{quote}

\section{Introduction}

This chapter reports the empirical evaluation of RepairA11y, the automated repair
system for WCAG 2.4 focus-behavior violations introduced in the previous chapter.
The purpose of the chapter is to answer, with measured data, the central question
that decides whether the system is worth building, and to position the result against
the closest published repair systems while stating honestly what the evidence can and
cannot support.

The central research question is stated here in plain language.

\begin{quote}
\textbf{RQ1 (evidence ablation): does runtime evidence help a large language model
(LLM) repair focus-behavior violations better than static evidence alone?}
\end{quote}

The question is answered by holding everything constant and varying only the evidence
bundle across four graduated levels, E1 to E4. Static evidence (the HTML and a
screenshot) is the lower end; runtime evidence (the computed styles and the measured
contrast ratio captured while the page is focused) is added at E3. The fourth level
(E4) adds a visual crop on top of the runtime evidence and acts as a control, testing
whether the effect keeps growing or saturates once runtime evidence is present.

Alongside effectiveness, the chapter reports a safety check on the same runs: do the
LLM patches introduce new conformance failures or change the page in a visually
disruptive way? A repair that fixes one problem but breaks another is not useful, so
this is reported next to every effectiveness number rather than as a separate study.

An earlier effectiveness comparison of generators, a separate rule-based generator,
and a focus-order criterion (SC 2.4.3) are out of scope for this chapter, which
focuses on the LLM repair path and the evidence it is given.

The headline findings are as follows. Runtime evidence (level E3) produces a large
and statistically significant improvement over static evidence (level E1) for LLM
repair. Adding the WCAG technique text alone (level E2) does not help, and on this
corpus it hurts. Adding a visual crop on top of runtime evidence (level E4) gives no
significant further gain, which shows the effect saturates once runtime evidence is
present. The patches are safe: across the 216 ablation trials no patch introduced a
new failure. Finally, the production scan shows that the target violation type is
common on real websites, which is the practical motivation for the work.

\section{Experimental Setup}

This section gives the detail needed to reproduce the experiments: the machine, the
software versions, the datasets, and the fixed configuration of the system. The
system has no training phase, so there are no learned hyperparameters to report. In
their place, the fixed configuration knobs and thresholds are listed, because those
are the values a reader would need to reproduce the numbers.

\subsection{Hardware and software environment}

The experiments were run on a single Apple macOS workstation (Darwin kernel
25.2.0, arm64). All code is JavaScript executed by Node.js (engine requirement
\verb|>=20.6.0|), with one Python script used only to draw the figures in this
chapter. The relevant software versions are fixed by the project's
\texttt{package.json} and are listed in Table~\ref{tab:env}.

\begin{table}[h]
\centering
\caption{Software environment.}
\label{tab:env}
\begin{tabular}{ll}
\toprule
Component & Version / setting \\
\midrule
Node.js & $\geq 20.6.0$ (uses the built-in \texttt{-{}-env-file} flag) \\
Playwright (headless Chromium) & \verb|^1.58.0| \\
pixelmatch (visual diff) & \verb|^7.2.0| \\
pngjs & \verb|^7.0.0| \\
sharp (image crop for E4) & \verb|^0.35.2| \\
ajv (JSON-schema validation) & \verb|^8.18.0| \\
vitest (unit tests) & \verb|^2.1.0| \\
LLM provider & OpenRouter, free tier \\
LLM model & \texttt{openai/gpt-oss-120b:free} \\
Decoding temperature & 0 (deterministic) \\
Provider rate limit & 3 requests per minute (free tier) \\
\bottomrule
\end{tabular}
\end{table}

The detection engine is NavA11y, the system from the previous chapter, used here
unchanged as both the violation detector (Stage 1) and the verification oracle
(Stage 5). NavA11y drives a headless Chromium browser through Playwright, calls
\texttt{.focus()} on each interactive element, and records 22 computed CSS properties
before and after focus, the measured contrast ratio of the focus indicator, and
whether the focused element is obscured by another element.

\subsection{The repair pipeline and its fixed configuration}

RepairA11y is a five-stage pipeline: detection, evidence packaging, generation, patch
application, and verification. The configuration knobs that affect the results, with
their fixed values, are listed in Table~\ref{tab:config}. These values were held
constant across every experiment in this chapter.

\begin{table}[h]
\centering
\caption{Fixed configuration (the system has no trained parameters).}
\label{tab:config}
\begin{tabular}{lll}
\toprule
Knob & Value & Meaning \\
\midrule
\texttt{temperature} & 0 & Deterministic LLM decoding, for reproducibility \\
\texttt{maxAttempts} & 3 & Retries when the LLM returns malformed JSON \\
\texttt{maxIterations} & 1 & Repair attempts per violation in the experiments \\
\texttt{ssimThreshold} & 0.9 & Visual-similarity floor for a regression flag \\
\texttt{RPM} & 3 & Free-tier request cap, a \texttt{\char`~}22\,s gate between calls \\
\bottomrule
\end{tabular}
\end{table}

The evidence packager is the experimental lever for RQ1. It produces one of four
graduated, strictly additive bundles. Each higher level contains everything in the
level below it, so the only difference between two adjacent levels is the new
information added.

\begin{itemize}
  \item E1: the element's outer HTML plus a full-page screenshot.
  \item E2: E1 plus the relevant WCAG technique text (for example, techniques F78,
    C27, G1).
  \item E3: E2 plus the runtime slice, namely the before-and-after computed style
    snapshots, the measured contrast ratio, and the obscurer element data.
  \item E4: E3 plus an annotated, cropped screenshot of the element with its bounding
    box drawn on.
\end{itemize}

\subsection{Datasets}

Two datasets are used. Their sizes and sources are summarized in
Table~\ref{tab:data}. The evidence ablation (RQ1), and the safety check measured on
the same runs, use the controlled corpus \Dd, where the ground truth is known and the
same cases can be repaired many times. The production corpus \Dr{} is used to measure
how common the target violation is in the wild and to test live repair.

\begin{table}[h]
\centering
\caption{Datasets. FAIL-case counts are per Success Criterion, as registered in
\texttt{src/datasets/}.}
\label{tab:data}
\begin{tabular}{p{1.6cm}p{2.6cm}p{5.5cm}p{3cm}}
\toprule
Dataset & Role & Size & Source \\
\midrule
\Dd & Controlled fixtures & 15 HTML files; FAIL cases: SC 2.4.11 = 2, SC 2.4.12 = 2, SC 2.4.13 = 6 & NavA11y focus-behavior fixture suite \\
\Dr & Production websites & 27 evaluable sites & Top-30 Semrush ranking, minus 3 exclusions \\
\bottomrule
\end{tabular}
\end{table}

Three of the top-30 production sites were excluded for access reasons, not for
convenience: Craigslist (geographic IP redirect), DoorDash (a bot-detection popup),
and MakeMyTrip (HTTP 403). The remaining 27 sites are the same corpus used by the
GenA11y detection study, which supports comparison.

One detail of the controlled dataset shapes the tables that follow. SC 2.4.7 has zero
FAIL cases in \Dd. This is not an omission. NavA11y classifies a completely absent or
invisible focus indicator under SC 2.4.13 (Focus Appearance), so cases that a reader
might expect under SC 2.4.7 (Focus Visible) appear under SC 2.4.13 instead. The
evidence ablation therefore runs on SC 2.4.13, which is where the controlled cases
concentrate and where the production violations are most common.

\subsection{Repeated trials and determinism}

The LLM generator is stochastic in principle, although temperature is pinned to 0 to
reduce variation. To measure the residual variation, the RQ1 experiment repeats each
cell: three independent runs over three seeds, so each evidence level is measured over
54 trials (6 cases times 3 seeds times 3 runs) rather than once.

\section{Evaluation Metrics}

This section defines every metric precisely, including the terms a reader could
misread. The most important definition is what counts as a successful repair.

\subsection{What counts as a successful repair}

A repair attempt on a single violation ends in exactly one status, produced by the
three-pass verifier in \texttt{src/verifier/index.js}.

\begin{itemize}
  \item \textbf{RESOLVED.} The target violation no longer fails when NavA11y is
    re-run on the patched page, \emph{and} the patch introduced no new failing
    element, \emph{and} the visual similarity between the page before and after the
    patch is at least the threshold of 0.9. All three conditions must hold.
  \item \textbf{UNRESOLVED.} The patch applied cleanly but the target still fails.
  \item \textbf{REGRESSED.} The patch introduced at least one new failing element, or
    it dropped visual similarity below 0.9.
  \item \textbf{DECLINED.} The generator explicitly returned no patch.
  \item \textbf{ERROR.} The attempt could not complete, for example because of an LLM
    rate-limit timeout, a malformed response after all retries, or a crash.
\end{itemize}

The primary metric is the resolution rate, defined in Equation~\ref{eq:res}.

\begin{equation}
\label{eq:res}
\text{resolution rate} = \frac{\text{number of RESOLVED trials}}{\text{total trials}}
\end{equation}

In the aggregate reporting code (\texttt{src/reporting/aggregate.js}), ERROR trials
are counted in the denominator. This is the conservative choice. Where errors are a
large share of trials, this chapter reports the resolution rate both with errors
included (the script default) and with errors excluded, and labels which is which.

\subsection{Regression rate}

The regression rate (Equation~\ref{eq:reg}) is the share of trials whose patch
introduced at least one new failure. It is the key safety metric, reported on the RQ1
runs.

\begin{equation}
\label{eq:reg}
\text{regression rate} = \frac{\text{trials with} \ge 1 \text{ new failing element}}{\text{total trials}}
\end{equation}

\subsection{Visual stability}

Visual stability measures how much the patch changed the rendered page, reported as a
similarity in $[0,1]$ where 1 means identical. Although the code calls this value
``SSIM'', the implementation in \texttt{src/verifier/ssim.js} is not structural
similarity; it is a per-pixel difference computed with the pixelmatch library
(Equation~\ref{eq:ssim}). This chapter uses the term ``pixel similarity'' to avoid
implying a computation that does not happen.

\begin{equation}
\label{eq:ssim}
\text{pixel similarity} = 1 - \frac{\text{mismatched pixels}}{\text{compared pixels}}
\end{equation}

When the before and after screenshots differ in size (which happens on live pages
that reflow), the comparison crops both to their common top-left overlap and reports
the fraction compared.

\subsection{Statistics for the evidence ablation}

RQ1 compares evidence levels on the same cases, so the outcomes are paired. The
correct test for paired binary outcomes is McNemar's test, not a $t$-test. The
implementation (\texttt{experiments/\_common/stats.js}) uses the continuity-corrected
form in Equation~\ref{eq:mcnemar}, with one degree of freedom, and reports a result
as significant when $p < 0.05$.

\begin{equation}
\label{eq:mcnemar}
\chi^2 = \frac{(\lvert b - c \rvert - 1)^2}{b + c}
\end{equation}

Here $b$ is the number of cases that succeeded at the lower evidence level but failed
at the higher one, and $c$ is the number that failed at the lower level but succeeded
at the higher one. Only discordant pairs enter the test.

The effect size between two proportions is Cohen's $h$ (Equation~\ref{eq:cohen}). The
conventional reading is that $\lvert h \rvert$ below 0.2 is negligible, 0.2 to 0.5
small, 0.5 to 0.8 medium, and 0.8 or above large.

\begin{equation}
\label{eq:cohen}
h = \left\lvert\, 2\arcsin\!\sqrt{p_1} - 2\arcsin\!\sqrt{p_2} \,\right\rvert
\end{equation}

For every reported proportion, a 95\% confidence interval is computed with the Wilson
score method, which is appropriate for small samples and for proportions near 0 or 1.

\subsection{Metrics that do not apply}

Several standard metrics are deliberately not reported, because the system does not
produce the quantities they need.

\begin{itemize}
  \item \textbf{Precision, recall, F1, accuracy:} not applicable. These are detection
    metrics; detection is the subject of the previous chapter. RepairA11y is a repair
    system, whose unit of success is a resolved violation.
  \item \textbf{Confusion matrix:} not applicable, for the same reason.
  \item \textbf{ROC curve and AUC:} not applicable. The system emits a typed patch,
    not a probability score.
  \item \textbf{Learning curve:} not applicable. There is no training phase; the LLM
    is used zero-shot.
\end{itemize}

\section{Experimental Results}

\subsection{Quantitative results}

\subsubsection{RQ1: evidence ablation (the central result)}

RQ1 holds everything constant and varies only the evidence bundle.
Table~\ref{tab:rq2dd} reports the per-level resolution rate on the controlled corpus,
and Table~\ref{tab:mcdd} reports the McNemar paired tests. Each level is measured over
54 trials (6 cases times 3 seeds times 3 runs).

\begin{table}[h]
\centering
\caption{RQ1 resolution rate per evidence level, \Dd{} (SC 2.4.13). Recorded run
\texttt{run-2026-06-21T06-19-00-943Z}.}
\label{tab:rq2dd}
\begin{tabular}{lrr}
\toprule
Level & Resolved / total & Rate \\
\midrule
E1 (static) & 36 / 54 & \SI{66.7}{\percent} \\
E2 (+ WCAG text) & 29 / 54 & \SI{53.7}{\percent} \\
E3 (+ runtime) & 50 / 54 & \SI{92.6}{\percent} \\
E4 (+ visual crop) & 48 / 54 & \SI{88.9}{\percent} \\
\bottomrule
\end{tabular}
\end{table}

\begin{table}[h]
\centering
\caption{RQ1 McNemar paired tests, \Dd. ``Significant'' means $p < 0.05$.}
\label{tab:mcdd}
\begin{tabular}{lrrrrrrl}
\toprule
Comparison & rate A & rate B & $b$ & $c$ & $\chi^2$ & $p$ & Cohen's $h$ \\
\midrule
E1 vs E2 & \SI{66.7}{\percent} & \SI{53.7}{\percent} & 7 & 0 & 5.143 & 0.0233 & 0.266 \\
E1 vs E3 & \SI{66.7}{\percent} & \SI{92.6}{\percent} & 2 & 16 & 9.389 & 0.0022 & 0.680 \\
E1 vs E4 & \SI{66.7}{\percent} & \SI{88.9}{\percent} & 2 & 14 & 7.563 & 0.0060 & 0.551 \\
E2 vs E3 & \SI{53.7}{\percent} & \SI{92.6}{\percent} & 2 & 23 & 16.000 & 0.0001 & 0.945 \\
E3 vs E4 & \SI{92.6}{\percent} & \SI{88.9}{\percent} & 5 & 3 & 0.125 & 0.7237 & 0.128 \\
\bottomrule
\end{tabular}
\end{table}

Three results stand out. First, moving from static evidence (E1) to runtime evidence
(E3) produces a large, significant gain of 25.9 percentage points ($p = 0.0022$, with
a medium-to-large effect size of $h = 0.680$). Second, adding the WCAG technique text
alone (E2) does not help and actually hurts, a significant drop of 13 points ($p =
0.0233$). Third, adding the visual crop on top of runtime evidence (E4) does not
significantly change the result ($p = 0.7237$, $h = 0.128$, negligible): the effect
has saturated once runtime evidence is present. Runtime evidence is the ingredient
that matters.

\subsubsection{Worked examples of LLM repairs on \Dd}

To make the repairs concrete, Table~\ref{tab:examples} lists the actual patch the LLM
produced at level E3 for each of the six controlled SC 2.4.13 cases, taken directly
from the recorded run (not reconstructed). Every patch is a CSS injection that adds a
compliant focus outline with \texttt{!important}, citing technique C15. Two details
are worth noting. First, the model chooses the colour to fit the element: it picks
black (\texttt{\#000}) for the dark-text buttons on a white background, where black
gives the most contrast, and a blue (\texttt{\#005fcc}) for the links, which matches
their existing styling while still meeting the contrast threshold. Second, it adds
\texttt{outline-offset} only where the element sits flush against other content, which
is the behavior runtime evidence enables: the model is reacting to the measured
layout, not guessing.

\begin{table}[h]
\centering
\caption{The LLM patch (level E3) for each \Dd{} SC 2.4.13 case, from the recorded
run \texttt{run-2026-06-21T06-19-00-943Z}.}
\label{tab:examples}
\small
\begin{tabular}{p{2.4cm}p{3.3cm}p{6.0cm}c}
\toprule
Case & What was wrong & LLM patch on the focused element (the fix) & Tech. \\
\midrule
No focus indicator & Focus produced no visible indicator & \texttt{outline: 2px solid \#005fcc !important; outline-offset: 2px} & C15 \\
Focus not indicated & Keyboard focus not shown visually & \texttt{outline: 2px solid \#005fcc !important; outline-offset: 2px} & C15 \\
Insufficient outline contrast & 2px outline at ${\sim}2.85{:}1$, below 3:1 & \texttt{outline: 2px solid \#000 !important; outline-offset: 0} & C15 \\
Insufficient outline width & Outline thinner than 2px & \texttt{outline: 2px solid \#000 !important} & C15 \\
Insufficient border width & Border-based indicator too thin & \texttt{outline: 2px solid \#000 !important} & C15 \\
Low background contrast & Indicator low-contrast against background & \texttt{outline: 2px solid \#000 !important} & C15 \\
\bottomrule
\end{tabular}
\end{table}

Figure~\ref{fig:examples} shows the focused element before and after the patch for
three of these cases. The change is deliberately small (a focus outline appearing or
becoming compliant), because that is exactly the repair the criterion requires; the
patch table above states precisely what changed in each case.

\begin{figure}[h]
\centering
\includegraphics[width=0.40\textwidth]{dd1_before.png}\hfill\includegraphics[width=0.40\textwidth]{dd1_after.png}\\[8pt]
\includegraphics[width=0.40\textwidth]{dd2_before.png}\hfill\includegraphics[width=0.40\textwidth]{dd2_after.png}\\[8pt]
\includegraphics[width=0.40\textwidth]{dd3_before.png}\hfill\includegraphics[width=0.40\textwidth]{dd3_after.png}
\caption{Three \Dd{} button cases, each shown focused before (left) and after (right)
the recorded LLM patch is applied. Top: a low-contrast 2px outline becomes a compliant
black one. Middle: a thin (1px) outline is thickened to a 2px black one. Bottom: a
thin border-based indicator is replaced by a 2px black outline.}
\label{fig:examples}
\end{figure}

\subsubsection{Regression and safety of the LLM patches}

Safety is measured on the same RQ1 runs, so it covers exactly the patches whose
effectiveness is reported above, with no separate experiment. Across all 216 ablation
trials the verifier flagged no regression: no patch introduced a new failing element,
and the mean pixel similarity between the page before and after the patch was 1.000.
Table~\ref{tab:safety} summarizes this.

\begin{table}[h]
\centering
\caption{Regression and safety of the LLM patches, measured on the RQ1 \Dd{} runs
(216 trials across the four evidence levels). Recorded run
\texttt{run-2026-06-21T06-19-00-943Z}.}
\label{tab:safety}
\begin{tabular}{lr}
\toprule
Metric & Value \\
\midrule
Trials evaluated & 216 \\
Trials that introduced a new failure & 0 (\SI{0.0}{\percent}) \\
Mean pixel similarity & 1.000 \\
\bottomrule
\end{tabular}
\end{table}

Because a patch that breaks the page is recorded as REGRESSED and never as RESOLVED,
the resolution rates in Table~\ref{tab:rq2dd} already exclude any unsafe patch. The
zero regression count means the effectiveness numbers are not bought at the cost of
breaking the page.

\subsubsection{\Dr: scale of the problem and live repair}

The production scan establishes that focus-behavior violations are common on real
sites. Table~\ref{tab:drdetect} reports the detection scan across 27 sites.

\begin{table}[h]
\centering
\caption{\Dr{} detection scan (recorded run; 26 of 27 sites scanned, 1 error).}
\label{tab:drdetect}
\begin{tabular}{lrrr}
\toprule
SC & Sites affected & Total violations & Mean per affected site \\
\midrule
2.4.7 & 18 & 270 & 10.38 \\
2.4.11 & 14 & 283 & 10.88 \\
2.4.12 & 21 & 497 & 19.12 \\
2.4.13 & 25 & 1841 & 70.81 \\
\bottomrule
\end{tabular}
\end{table}

SC 2.4.13 dominates: 25 of 26 successfully scanned sites have at least one Focus
Appearance violation, with 1{,}841 in total. Live repair on production sites is the
hardest test. Two numbers are reported with equal weight, because they answer
different questions (Table~\ref{tab:drrepair}).

\begin{table}[h]
\centering
\caption{\Dr{} live repair on production sites (SC 2.4.13, level E3, recorded runs).}
\label{tab:drrepair}
\begin{tabular}{p{2.6cm}p{6.5cm}r}
\toprule
Reading & Definition & Result \\
\midrule
Single-run rate & One attempt per site, errors excluded & 10 / 21 = \SI{48}{\percent} \\
Best-of rate & Best outcome per site across retries & 18 / 25 = \SI{72}{\percent} \\
\bottomrule
\end{tabular}
\end{table}

The single-run rate (\SI{48}{\percent}) is the honest per-attempt expectation on a
noisy free tier: of 25 sites, one run resolved 10, left 9 unresolved, regressed 2,
and hit 4 rate-limit or timeout errors. The best-of rate (\SI{72}{\percent}) is the
upper bound when each site is retried, and it should be read with two caveats: it
counted one site (adp.com) as a regression from a dry-run stub with no real API call,
and one site (live.com) regressed badly, with a pixel similarity of 0.493 and 60 new
failures.

\subsection{Visual results}

Figure~\ref{fig:rq2dd} shows the RQ1 evidence ablation on the controlled corpus: the
jump at E3 and the flat step to E4 are both visible. Figure~\ref{fig:drsc} shows the
scale of violations across production sites.

\begin{figure}[h]
\centering
\includegraphics[width=0.62\textwidth]{rq2_ablation_dd.png}
\caption{RQ1 evidence ablation on the controlled corpus \Dd. The rate rises sharply
from E1 to E3 and is flat from E3 to E4.}
\label{fig:rq2dd}
\end{figure}

\begin{figure}[h]
\centering
\includegraphics[width=0.62\textwidth]{dr_violations_by_sc.png}
\caption{Total focus-behavior violations detected across production sites, by SC.}
\label{fig:drsc}
\end{figure}

Before-and-after screenshots of the controlled cases are shown in
Figure~\ref{fig:examples}, using the real recorded LLM patches. Screenshots of live
production repairs are not shown: the only production captures available were not
genuine repairs (one site was a dry-run stub with no patch applied, another returned a
bot-challenge page rather than the real site), so the production results are reported
as numbers (Table~\ref{tab:drrepair}) rather than as images. No machine-learning
training figures (learning curves, ROC curves) are shown either, because the system
has no training phase and emits no probability score; those figures are not
applicable here.

\section{Comparative Analysis}

\subsection{Against a no-op baseline}

The simplest baseline is to do nothing, which by construction resolves
\SI{0}{\percent} of violations. Every result in this chapter is therefore measured
against a \SI{0}{\percent} floor.

\subsection{Static versus runtime evidence}

The cleanest internal comparison is E1 against E3, because the only difference between
them is the runtime slice. Runtime evidence lifts the LLM resolution rate by 25.9
points, a statistically significant change. This is reported as a controlled ablation
rather than a comparison against an external tool, because no external tool exposes
the same lever.

\subsection{Against published repair systems}

Table~\ref{tab:related} places RepairA11y next to the closest published LLM repair
systems. The comparison must be read carefully: these systems target different Success
Criteria, use different detectors, and report different metrics. The numbers for prior
systems are taken from their papers, cited in \texttt{docs/LITERATURE.md}.

\begin{table}[h]
\centering
\caption{Positioning against published work. ``Ablation'' is whether the system
measures the contribution of different evidence inputs.}
\label{tab:related}
\begin{tabular}{p{3.2cm}p{1.8cm}p{2.2cm}p{2.6cm}cc}
\toprule
System & Task & Detector & Reported result & Ablation & Focus SCs \\
\midrule
GenA11y (FSE 2025) & Detection & GPT-4o + DOM & \SI{94.5}{\percent} P, \SI{87.6}{\percent} R & no & partial \\
AccessGuru (ASSETS 2025) & Repair & axe + LLM & up to \SI{84}{\percent} score reduction & no & no \\
Fern\'andez-Navarro \& Chicano (2026) & Repair & axe + Selenium & \SI{80}{\percent} static, \SI{86}{\percent} Angular & no & no \\
DesignRepair (ICSE 2025) & Frontend & source + view & \SI{92.9}{\percent} R, \SI{90.1}{\percent} P & no & no \\
RepairA11y (this work) & Repair & NavA11y runtime & E3 \SI{92.6}{\percent} on \Dd & yes & yes \\
\bottomrule
\end{tabular}
\end{table}

The published repair systems report fix rates in the \SI{80}{\percent} to
\SI{86}{\percent} range, broadly comparable to RepairA11y's runtime-evidence result.
The contribution of this chapter is not a higher headline number on a shared
benchmark, which does not exist for these SCs. It is the evidence ablation: none of
the prior systems measures whether runtime evidence helps, because each uses a single
fixed evidence model, and most are bounded by what the axe static analyzer reports.

\section{Statistical Analysis}

\subsection{Choice of test}

The RQ1 comparisons are paired: the same cases are repaired at each evidence level.
For paired binary outcomes the appropriate test is McNemar's test. A two-sample
$t$-test would be wrong, because the outcomes are binary and paired, not continuous
and independent. A chi-square test of independence would also be wrong, because it
ignores the pairing. McNemar is reported with the continuity correction
(Equation~\ref{eq:mcnemar}) and one degree of freedom. For the smallest discordant
counts the chi-square approximation is weaker than an exact binomial test; this is
noted as a limitation in Section~\ref{sec:threats} rather than corrected, because the
significant results have discordant counts large enough for the approximation to hold.

\subsection{Confidence intervals}

Because the proportions are estimated from a small sample (6 SC 2.4.13 cases, repeated
to 54 trials per level), Wilson score 95\% intervals are reported instead of
normal-approximation intervals. Table~\ref{tab:wilson} reports the interval for each
per-level proportion.

\begin{table}[h]
\centering
\caption{Wilson 95\% confidence intervals for the RQ1 per-level rates (recorded run).
$n$ is the number of trials behind each proportion.}
\label{tab:wilson}
\begin{tabular}{lrrl}
\toprule
Level & Rate & $n$ & Wilson 95\% CI \\
\midrule
E1 & \SI{66.7}{\percent} & 54 & (\SI{53.4}{\percent}, \SI{77.8}{\percent}) \\
E2 & \SI{53.7}{\percent} & 54 & (\SI{40.6}{\percent}, \SI{66.3}{\percent}) \\
E3 & \SI{92.6}{\percent} & 54 & (\SI{82.4}{\percent}, \SI{97.1}{\percent}) \\
E4 & \SI{88.9}{\percent} & 54 & (\SI{77.8}{\percent}, \SI{94.8}{\percent}) \\
\bottomrule
\end{tabular}
\end{table}

The intervals for E1 and E3 do not overlap, which is consistent with the significant
McNemar result. The E3 and E4 intervals overlap heavily, consistent with the
non-significant step from E3 to E4.

\subsection{Effect sizes}

The effect of runtime evidence is not only statistically significant, it is large.
Cohen's $h$ for E1 versus E3 is 0.680 (medium to large). By contrast, E3 versus E4 has
$h = 0.128$, which is negligible.

\section{Resource Utilization}

The monetary cost is zero. Every LLM call uses the OpenRouter free tier, so no
experiment in this chapter incurred a charge. The free tier is the reason the project
pins temperature to 0 and uses a single free model.

The dominant time cost is the free-tier rate limit, not computation. The provider
allows three requests per minute, which the client enforces as a gate of about 22
seconds between calls. On top of this, each trial runs NavA11y twice (baseline
detection and post-patch verification), and each pass drives a real headless browser,
taking roughly 15 to 30 seconds. A single LLM trial therefore takes on the order of a
minute, and RQ1, at 216 trials across the four levels, takes hours of wall-clock time.
The repair loop used a mean of 1.0 iterations.

Token counts and per-call latency are not written to the result files in a form this
chapter can total reliably. Reporting an exact token total is therefore marked
\texttt{[VERIFY]}: confirming it would require adding token accounting to the LLM
client and re-running, which was out of scope for a zero-budget evaluation.

\section{Discussion of Findings}

The central finding is that runtime evidence makes LLM repair work, and the reason is
visible in the failure cases. At levels E1 and E2 the model can see the HTML and, at
E2, the WCAG technique text, but it cannot see the computed result of the page's CSS.
For a Focus Appearance violation the deciding facts are the actual rendered outline
width, the actual indicator colour, and the actual background colour behind it, none
of which can be read reliably from the source because they come from external
stylesheets, CSS resets, and custom properties. Without those facts the model guesses,
and its most common guess is a black outline, which fails on dark backgrounds. At E3
the model is given the measured contrast ratio and the before-and-after computed
styles, so it stops guessing. This directly answers RQ1.

The result that E2 does not help, and on this corpus hurts, is consistent with the
literature that more context can degrade an LLM. Adding the WCAG text without the
runtime numbers gives the model more to read but nothing new to reason with. This is
why the chapter does not treat ``E3 is better than E1'' as obvious: a plausible
competing hypothesis, that any extra context helps, is contradicted by E2.

The result that E4 does not beat E3 is the reason E4 was included at all. E4 is the
control that tests whether the gain keeps growing once runtime evidence is present. It
does not: the annotated crop adds visual grounding, but once the model already has the
measured numbers, the picture is largely redundant. The practical reading is that
runtime evidence is the sweet spot, and the system does not need to pay the extra cost
of rendering and cropping an annotated image. This answers the natural follow-up
question, ``would a picture help even more?'', with measured data: no.

On safety, the zero regression rate and the pixel similarity of 1.000 across the 216
ablation trials say the LLM patches do not break the controlled pages: the
effectiveness gain is not bought by damaging the rest of the page. The honest
qualifier is that this is demonstrated on fixtures, and the production runs show a
regression tail (notably live.com) that the controlled suite does not capture. Safety
on fixtures is necessary but not sufficient for safety in the wild.

The drop from controlled to production performance (a single-run \SI{48}{\percent} on
\Dr) is explained by three compounding factors: the free tier loses roughly one trial
in fourteen to errors, the oracle's transparent-background bug blocks some correct
patches, and live pages are larger and messier than fixtures. None of these undermines
the RQ1 finding, which is an internal comparison where all three factors apply equally
to every level.

\section{Threats to Validity and Limitations}
\label{sec:threats}

\textbf{Small, single, in-house corpus.} The evidence ablation rests on one
controlled corpus of 6 SC 2.4.13 cases (repeated to 54 trials per level), authored by
the same project. This is the chapter's main external-validity limit: the controlled
effect is measured cleanly, but its generalization to pages the team did not write is
argued, not yet proven on a held-out repair benchmark. The production scan (\Dr)
shows the problem is real at scale, and the production repair run gives a first
out-of-lab data point, but the controlled ablation itself is in-house.

\textbf{Wide confidence intervals.} With 54 trials per level the Wilson intervals are
wide (Table~\ref{tab:wilson}). The significant RQ1 results survive this because they
are paired and the effect is large, but the per-level point estimates should be read
as indicative.

\textbf{Repeated-case pairing in McNemar.} The McNemar test is applied to pairs formed
by (case, run, seed), which treats repeats of the same case as independent trials.
This inflates the effective sample size and narrows the p-values relative to a test
that pairs only at the case level. The direction and size of the effect are not in
doubt, but the exact p-values should be read with this in mind.

\textbf{LLM stochasticity.} Temperature 0 and repeated runs reduce but do not
eliminate variation. A fresh run is a new sample and will not reproduce the recorded
numbers exactly.

\textbf{Oracle conservativeness.} The verifier is NavA11y itself, which for elements
with a transparent background treats the background as black when computing contrast,
so a correct outline patch can be marked as still failing. This makes the reported
rates conservative for those cases.

\textbf{Free-tier flakiness.} Roughly \SI{7}{\percent} of trials are lost to
rate-limit timeouts, malformed JSON after retries, or browser crashes. These are
counted as ERROR and, by default, against the resolution rate, making the headline
rates a lower bound.

\textbf{Optimistic best-of reporting.} The \SI{72}{\percent} best-of figure on
production sites is an upper bound obtained by retrying each site, and it includes at
least one site counted from a dry-run stub. It is reported alongside the single-run
\SI{48}{\percent} so it is not mistaken for the expected per-attempt result.

\textbf{Single Success Criterion for the ablation.} RQ1 is run on SC 2.4.13 only. The
ablation finding is established for Focus Appearance; generalization to the other SCs
is plausible but not yet measured.

\textbf{Oracle overfitting.} Because the generator is verified by the same engine that
detects, there is a risk of fitting the oracle rather than fixing the page. This is
the largest reviewer concern. The chapter reports the oracle pass rate; a full
independent manual review remains future work.

\section{Summary}

This chapter evaluated RepairA11y on a controlled corpus and on production data around
one central research question. The central result, for RQ1, is that runtime evidence
(E3) improved LLM repair over static evidence (E1) by 25.9 points ($p = 0.0022$,
medium to large effect), while adding WCAG text alone (E2) gave no gain and hurt, and
adding a visual crop (E4) gave no significant further gain, showing the effect
saturates at runtime evidence. The same patches were safe: across the 216 ablation
trials none introduced a new failure, with mean pixel similarity of 1.000. On
production sites, focus-appearance violations are common (1{,}841 across 25 sites),
one-shot repair resolved about half of the targeted violations, and retried repair
resolved about three-quarters, with a regression tail that must be watched.

The practical message is that the value of an LLM accessibility repair system comes
less from the model and more from the runtime evidence it is given, and that this
contribution can be measured. The next chapter draws the thesis together and discusses
how runtime-grounded repair could be extended to the remaining focus-behavior Success
Criteria and hardened for production use.

\end{document}
